% eksperymentalne tłumaczenie na włoski

%

Konstrukcje od \ref{delta_2024_12_start} do \ref{delta_2024_12_end} opisze Stanisław Majchrzak w $\Delta_{24}^{12}$. % lang-pl
Le costruzioni da \ref{delta_2024_12_start} a \ref{delta_2024_12_end} sono descritte da Stanisław Majchrzak in $\Delta_{24}^{12}$. % lang-it

\begin{geoconstruction}
\label{delta_2024_12_start}%
    Dane jest pięć punktów leżących na okręgu $\Gamma$. % lang-pl
    Wykreślić styczną do $\Gamma$ w jednym z tych punktów. % lang-pl
    %
    Sono dati cinque punti appartenenti a un cerchio $\Gamma$. % lang-it
    Costruire la tangente a $\Gamma$ in uno di questi punti. % lang-it
\end{geoconstruction}

Uzasadnienie poprawności tej konstrukcji wymaga znajomości twierdzenia Pascala, patrz $\Delta_{14}^9$. % lang-pl
\index{twierdzenie!Pascala}% % lang-pl
La dimostrazione della correttezza di questa costruzione richiede la conoscenza del teorema di Pascal; si veda $\Delta_{14}^9$. % lang-it
\index{teorema!di Pascal} % lang-it

\begin{geoconstruction}
    Dane jest pięć punktów leżących na okręgu $\Gamma$ oraz prosta $l$ przechodząca przez jeden z nich. % lang-pl
    Wyznaczyć drugi punkt przecięcia okręgu $\Gamma$ oraz prostej $l$. % lang-pl
    %
    Sono dati cinque punti appartenenti a un cerchio $\Gamma$ e una retta $l$ passante per uno di essi. % lang-it
    Determinare il secondo punto d'intersezione tra il cerchio $\Gamma$ e la retta $l$. % lang-it
\end{geoconstruction}

Rozwiązanie zostanie wydrukowane w $\Delta_{17}^6$. % lang-pl
La soluzione sarà pubblicata in $\Delta_{17}^6$. % lang-it

\begin{geoconstruction}
    Dane są dwa okręgi przecinające się w dwóch punktach. % lang-pl
    Wyznaczyć środek jednego z nich. % lang-pl
    %
    Sono dati due cerchi che si intersecano in due punti. % lang-it
    Determinare il centro di uno di essi. % lang-it
\end{geoconstruction}

W kolejnych konstrukcjach przydatna jest znajomość twierdzenia: biegunowe dowolnego punktu $P$ względem okręgów należących do jednego pęku są współpękowe. % lang-pl
Nelle costruzioni successive è utile conoscere il seguente teorema: le polari di un punto arbitrario $P$ rispetto ai cerchi appartenenti a uno stesso fascio sono concorrenti. % lang-it

\begin{geoconstruction}
    Dane są dwa okręgi oraz punkt $A$ na zewnątrz jednego z nich (wewnątrz obydwu). % lang-pl
    Wykreślić okrąg przechodzący przez $A$, który jest współpękowy z dwoma danymi. % lang-pl
    %
    Sono dati due cerchi e un punto $A$ esterno a uno di essi (e interno a entrambi). % lang-it
    Costruire un cerchio passante per $A$ che appartenga allo stesso fascio dei due cerchi assegnati. % lang-it
\end{geoconstruction}

\begin{geoconstruction}
\label{delta_2024_12_end}%
    Dane są cztery (w trudniejszej wersji: trzy) okręgi takie, że żadne trzy nie są współpękowe. % lang-pl
    Wyznaczyć środek przynajmniej jednego z nich. % lang-pl
    %
    Sono dati quattro cerchi (nella versione più difficile: tre) tali che nessuna terna appartenga allo stesso fascio. % lang-it
    Determinare il centro di almeno uno di essi. % lang-it
\end{geoconstruction}

Dla rozłącznych okręgów należących do jednego pęku, środka nie da się (!) skonstruować. % lang-pl
Per cerchi disgiunti appartenenti a uno stesso fascio, il centro non può (!) essere costruito. % lang-it

%